\chapter{Vectors}
\label{Vector_Chapter}

\section{Just an arrow}

\noindent The mass of an object or its temperature can be expressed by a number (with units): 10 kg, 37$^{\circ}C$, 23 pounds, 103$^{\circ}$Fahrenheit,...\\

\noindent There are quantities that cannot be described by only a number. If I push against the table to the right it will move to the right. But pushing it to the left will make it move to the left. Therefore, only knowing how hard I push against the table does not suffice to determine what will happen to the table. For this we need the \textbf{strength} as well as the \textbf{direction}. We therefore introduce arrows. The size or length of the arrow represent the strength of the force and the direction is determined by where the arrow points to. These arrows are called vectors and we will briefly discuss their properties and how we can use them to describe physical quantities where only a number does not suffice.\\

\noindent Below is a vector. As you can see it is just an arrow with some length and pointing in some direction. As naming objects is knowing them, we decide to call this one $\vec{a}$. When talking about vectors we will always write the arrow above the letter to emphasize the fact that is indeed a vector and not merely a number. Its size on the other hand is a (positive) number and we will denote it by $a$. So the vector $\vec{a}$ has size $a$.\\

\noindent A vector has 3 main properties:
\begin{enumerate}[a]
\item \textbf{Length: }the distance from its starting point to its end
\item \textbf{Direction: }the direction towards which the arrow points
\item \textbf{Starting point: }the point where it starts
\end{enumerate}
Two vectors will be considered equal if their size and direction are the same, even if we draw them at different points. We will ignore their starting point unless stated otherwise\footnote{The starting point will play a role when discussing rotational motion, angular momentum and torque.}. Therefore we are free to take any vector and move it around. As long as we do not rotate it or change its length, it remains the same vector.

\section{Vector arithmetic}

\noindent Now that we have introduced vectors we need to learn how to work with them. 

\subsection{Addition of vectors}

\noindent The sum of two vectors $\vec{a}$ and $\vec{b}$ is defined on the figure below:

\noindent The sum $\vec{a}+\vec{b}$ is itself a vector connecting the starting point of $\vec{a}$ with the end point of $\vec{b}$ when the latter is drawn so that that its starting point coincides with the end point of the former.

\subsection{Multiplication of a vector with a number}

\subsection{Subtraction of vectors}

\subsection{Multiplication of vectors}

\subsubsection{Dot product of vectors}

\subsubsection{Cross product of vectors}

\section{The components of vectors}

\section{The components of operations of vectors}

\section{Scalar and vector fields}

\section{Vector calculus}

\subsection{Del operator: gradient, divergence, curl and Laplacian}

\subsection{Green's theorem}

\subsection{Stokes' theorem}



%
%\section{Vector calculus}
%
%\section{Vector as an abstraction of arrows}